\section{Conclusions and evolution to Stage 3}
\label{sec:stage3}

\subsection{Strengths and weaknesses of the coordination}

\textbf{Strengths.} The coordination is explicit and checked: every transition is one method of one library, atomic,
and tested against real Hazelcast members. A book never changes state before its evidence is persistent, so a crash
at any point leaves a book that is either redone or recovered; the failure tests check it end to end
(Table~\ref{tab:failure-tests}). Replicas of
Ingestion and Search are added with \code{--scale} and nothing else, because nobody addresses a replica: the workers
pull from shared queues, and nginx spreads the queries. Losing the whole Hazelcast cluster costs a recovery, not data.

\textbf{Weaknesses.} Hazelcast is in memory, so the books without evidence (\code{PENDING}, or \code{INGESTING} without
\code{\_READY}) must be sent again after the cluster is lost. The lock of the index writer has no fencing token. The
leases compare the clocks of different processes. The datalake depends on a local file system and its atomic moves. Each
book costs two transactions, and both queues are single partitions that every claim goes through.

\pending{The measured part: books per second of Ingestion with 1, 2 and 4 replicas and with 1, 2, 4 and 8 workers, the
share of coordination in the time of a book, and the queue wait (benchmarks/ingestion/results/integration/), with the
bottleneck each one shows.}

\subsection{What changes on the Stage 3 cluster}

The services keep their code; what they assume about the machine does not hold on a cluster. Table~\ref{tab:stage3}
lists each local assumption and what replaces it.

\begin{table}[H]
\centering\small
\begin{tabularx}{\linewidth}{@{}L{3.6cm}L{4.3cm}X@{}}
\toprule
\textbf{Local assumption} & \textbf{Why it breaks on a cluster} & \textbf{Stage 3} \\
\midrule
Two Hazelcast members on one host & They share its failures; a backup on the same machine protects from nothing &
  One member per machine (or embedded in the services), discovery by the platform instead of a fixed list, and backups
  placed on another machine (partition groups per host or zone) \\
\code{IMap.lock} for the single writer & In a network partition each half of the cluster can grant it (split-brain), and
  it has no fencing token & A CP \code{FencedLock} (Enterprise, three or more CP members), or a lease document in MongoDB
  taken with a conditional update whose version is the fencing token for every write; a split-brain protection quorum
  on the coordination structures \\
Leases on one clock & Clocks of different machines drift; a lease could expire early or late & NTP on every node and a
  margin above the expected skew, or the time of the member that owns the entry \\
Datalake on a local file system, \code{\_READY} after atomic moves & An atomic move exists only inside one file system;
  Indexing on another machine cannot read the folder & Shared storage (an object store such as S3 or MinIO, or HDFS): write
  the book, then the marker object; readiness is the marker object \\
Recovery lists the datalake & Listing every range folder grows with the number of books, and gets slower over the network &
  A persistent log of ready books (a MongoDB collection written with the marker), read like \code{index\_log} \\
One Control & Its absence stops leases and publications; on a cluster, nodes come and go & Two or more instances with a
  leader election (the same fenced lease as the writer); the others wait as standbys \\
Queues of one partition & Every claim of every node goes to one member & Several queues by range of ids, each taken by
  some workers, or a partitioned log with acknowledgements (Kafka) if the work grows beyond what one member serves \\
Services without authentication on \code{127.0.0.1} & On a network of machines anyone could reach the Control or
  Hazelcast & TLS and authentication on the Hazelcast clients and on the internal APIs; only the load balancer exposed \\
Docker network on one host & Containers on different machines do not share it & An overlay network or an orchestrator
  (Kubernetes, Docker Swarm) with service discovery; nginx replaced by its load balancer \\
\bottomrule
\end{tabularx}
\caption{The local assumptions of the coordination and their Stage 3 replacement.}
\label{tab:stage3}
\end{table}

\textbf{What does not change.} The contract itself: the states, the transitions, the principle that evidence comes
first, the leases and the recovery are the same on one machine and on many, and the services reach Hazelcast through the
same library. That was the reason to choose it (Table~\ref{tab:alternatives}): for Stage 3 the work is in the
configuration and in the storage, not in the services.

\textbf{Distribution of work.} On a cluster the replicas of Ingestion run on several machines and keep pulling from the
same queues, so the distribution of books needs no change. What needs measuring again is where the first bottleneck moves
once each replica has a machine of its own: the queues of Hazelcast, the shared storage, or the single index writer.
\pending{Link with the conclusion of section 7: which of these limits the single-machine measurements already show.}
